\chapter{Conclusion}
In this thesis, a first-time analysis of the cH process in the \cHZZ channel is presented. To this end, the 2018 data taking era of the CMS detector is considered, with the ultimate goal of placing constraints on the charm quark Yukawa interaction modifier \kappac. Initially, the SM of particle physics is discussed with a particular focus on the Higgs boson, as it is the Higgs boson which is responsible for the generation of particle masses in the SM. This is also true for the charm quark, which gains its mass through a Yukawa interaction with the Higgs boson. Unlike the Yukawa interactions of other heavier fermions in the SM, the charm-Yukawa interaction is yet to be measured experimentally. Motivated by this, the cH process is identified as a process that is of experimental interest for such a measurement and is discussed in the context of the SM. A method by which a signal strength measurement of the cH process can be translated to a constraint on \kappac\, is also introduced. Subsequently, the interpretation of the cH process in a beyond-the-SM physics scenario, parametrised through SMEFT, is discussed. Here, specifically the chromomagnetic dipole operator is identified as a SMEFT operator that may be uniquely constrained with the cH process. This is followed by a discussion of the latest results in analyses targeting \kappac. \\
\\
Following this, the CMS detector is introduced. This is a multipurpose particle physics detector that can reconstruct pp collisions at the LHC to identify cH candidate events. The various subdetector systems of the CMS detector and the reconstruction methods and algorithms relevant to the analysis are discussed. This includes jet flavour tagging algorithm \textsf{DeepJet}. Additionally, the Monte Carlo simulation methods that can be used to estimate the physics processes that occur in the pp collisions recorded by the CMS detector at the LHC are presented. \\
\\
Subsequently, the analysis strategy employed in this work to target the \cHZZsen process is detailed. This includes the reconstruction of Higgs boson candidates and the selection of an associated jet candidate. From these individual components, candidate \cHZZ events are obtained. The simulation of the \cHZZsen signal process is discussed, alongside the methods that are used to estimate a variety of background processes. This is followed by a presentation of the statistical methods that are used to derive expected 95\% CL upper limit on the signal strength of the \cHZZsen process, as well as the model of systematic uncertainty sources that is implemented. With this method, expected 95\% CL upper limits on \sBR as well as \kappac\, are set for a number of scenarios. In the nominal scenario where the Higgs boson candidate mass is used as a discriminator, this corresponds to an expected 95\% CL upper limit of \sBR$= 1088$ times the SM expectation ($\mid$$\kappa_c$$\mid$ = 191) for the 2018 dataset recorded by the CMS detector. Additional scenarios using different discriminators are considered and, for the nominal scenario, a statistical extrapolation to a full Run-2 dataset and the anticipated High-Luminosity LHC dataset is performed. The corresponding expected 95\% CL upper limits that are derived are \sBR$ = 725$ times the SM expectation ($\mid$$\kappa_c$$\mid$ = 128) and \sBR$ = 318$ times the SM expectation ($\mid$$\kappa_c$$\mid$ = 58), respectively. Additionally, the impact of leading systematic uncertainty sources in the evaluation procedure are discussed, along with a fit to data in sidebands region of the analysis. Potential analysis improvements are also discussed, including the modelling uncertainties that affect analyses of the cH process. \\
\\
Finally, the potential of the presented analysis to constrain the Wilson coefficient of the chromomagnetic dipole operator in SMEFT is demonstrated and discussed. An expected constraint of $\tilde{C}_{cG} < 1.94$ is set at 95\% CL. Again, a statistical extrapolation to full Run-2 and High Luminosity-LHC scenarios is performed, with the expected constraints being set at $\tilde{C}_{cG} < 1.58$ and $\tilde{C}_{cG} < 0.75$ at 95\% CL, respectively. Potential improvements to this interpretation of the analysis are also discussed. \\
\\
The expected sensitivity to \kappac\, that is derived in this work is lower than in other, similar analyses. This includes analyses of the cH(WW) and cH($\gamma\gamma)$ processes, as well analyses targeting H$\rightarrow \mathrm{c\overline{c}}$ decays. However, the presented analysis remains to be an important, complementary contribution in the quest to determine the nature of the charm-Yukawa coupling, especially given the significant challenges associated with such a measurement. Furthermore, significant insights are gained with respect to the jet selection that must be performed in cH analyses in light of the kinematics-based jet association algorithm that is developed here. Additionally, key systematic uncertainties such as the flavour scheme uncertainty in modelling the cH process are discussed and potential solutions to this uncertainty are proposed. Finally, the novel results of the presented SMEFT interpretation of the cH process represent a unique method with which to probe certain beyond-the-SM scenarios that may provide exciting hints at deviations from expected SM behaviour. 
\section*{Outlook}
As data taking at the LHC draws to a close and particle physicists look towards the HL-LHC, the measurement of the charm-Yukawa coupling remains to be a highly relevant, yet challenging facet of our characterisation of the Higgs boson. A key component of this challenge arises from the very small cross sections associated with processes sensitive to this coupling. Accordingly, the continued development of a variety of analysis strategies that target such processes remains to be a key strategic component in ultimately making such a measurement. As a complement to analysis channels in which the Higgs boson decays directly to a pair of charm quarks, strategies targeting charm quark-associated Higgs production have gained increased attention. Given their novel status, the continued development of such analysis strategies remains to be of significant interest and is a multifaceted endeavour, considering the variety of Higgs boson decay channels that may be targeted. These include, so far, the H$\,\rightarrow4\mu$ channel presented in this work, as well as the H$\,\rightarrow WW$ and H$\,\rightarrow \gamma\gamma$ channels. However, channels such as H$\,\rightarrow 4e$, H$\,\rightarrow 2e2\mu$, and H$\,\rightarrow \tau\tau$ remain largely unexplored, much like SMEFT interpretations of the cH process. With a variety of improvements, such as developing experimental methods to specifically target low \pt \, jets or improvements in the simulation of the cH process, significant gains in the sensitivity of these analyses to the charm-Yukawa coupling or a SMEFT interpretation thereof may be achieved. These improvements are especially relevant given that, at first order, they are independent of the Higgs boson decay channel of the cH process. With the many different analysis strategies and improvements to be explored with the full Run-2, Run-3 and HL-LHC datasets, the future of a charm-Yukawa coupling measurement surely continues to look promising.